\section{研究ストーリーと仮説}
\label{sec:story}

\subsection{研究の流れ（ストーリーライン）}

本研究は，次の論理構造に基づいている．

\begin{enumerate}[leftmargin=2em]
  \item \textbf{問題}: 工場内でAGVと人が共存する際，衝突リスクの予測が困難
  \item \textbf{既存アプローチの限界}: 経路を常時・均一に表示するeHMIは，
        危険の程度を区別できず，注意の優先順位付けを支援しない
  \item \textbf{提案}: 残存経路を提示しつつ，
        \emph{相対的危険度}に応じて色相（連続）と不透明度（段階）を変化させる
  \item \textbf{評価}: VR上の工場タスク（Station A$\rightarrow$B移動）で，
        効率（時間・移動距離）と安全（ニアミス）を定量的に比較
  \item \textbf{期待}: Proposedは，No-ARより安全に，Baselineより情報過多を抑えつつ
        適切な注意喚起を実現する
\end{enumerate}

\subsection{理論的根拠}

危険度の時間的予測として\textbf{TTC（Time-to-Conflict）}を用いる点は，
交通・ロボット分野で広く用いられる指標に依拠する\cite{markkula2020}．
また，eHMI研究では，外部表示が歩行者の意図理解や信頼に影響することが報告されている\cite{francis2020eHMI,matthiesen2021survey}．

本提案の新規性は，次の3点に集約される．

\begin{itemize}[leftmargin=2em]
  \item \textbf{視線追従型の動的基準線}:
        目的地最短経路ではなく，参加者が\emph{現在注視している方向}に
        AGVが割り込みうる領域を定義する
  \item \textbf{二重危険度指標の統合}:
        交差予測（TTC）と近接距離の最大値により，
        「交差しないが近い」「交差するがまだ遠い」両ケースをカバー
  \item \textbf{連続色＋段階的不透明度}:
        危険の\emph{連続的変化}（色相）と\emph{警告レベル}（不透明度4段階）を
        分離してエンコード
\end{itemize}

\subsection{研究仮説}

被験者内計画に基づき，以下の仮説を設定する．

\begin{description}[leftmargin=0em,style=nextline]
  \item[\textbf{H1（安全性）}]
    Proposed条件は，No-AR条件と比較して\textbf{ニアミス回数が少ない}．
    根拠: 危険AGVの経路が視覚的に強調され，回避行動が促進される．

  \item[\textbf{H2（効率）}]
    Proposed条件は，Baseline条件と比較して\textbf{完了時間が短い}，
    または\textbf{移動距離が短い}．
    根拠: 危険度に応じた選択的強調により，不要な情報処理負荷が低減される．

  \item[\textbf{H3（情報なし対照）}]
    No-AR条件は，Proposed条件と比較して\textbf{完了時間が長い}，
    または\textbf{移動距離が長い}．
    根拠: 経路情報がないと，AGVの動き予測が困難になり迂回・停止が増える．
\end{description}

\noindent
\textbf{注記:}
H1--H3は独立ではなく，ProposedがBaselineとNo-ARの\emph{中間的な最適点}を
示すことを期待する．
統計検定は，被験者内反復測定のANOVA（または Friedman 検定）と
事後比較（Bonferroni 補正）を想定する．
